% Transcription — fonds Alexandre Grothendieck, shelfmark 88, batch 1 (pages 1-15).
% Established from the facsimile archives/batches/88/batch-01.pdf (a mirror of
% https://grothendieck.umontpellier.fr/88.pdf), per the skill
% transcribe-grothendieck. The folder is fifteen pages and this is its only
% batch.
%
% Thirteen of the fifteen leaves are transcribed, and the folder holds two runs
% rather than one. Page 1 is a tan wrapper carrying, alone in the top right
% corner and in pencil, « Cartes associées aux syst. de Racines » followed by a
% « (7) »; it carries nothing else, so it takes no \page marker, its words
% become the first section heading, and they are his and not the archivists'.
% Page 9 is the other face of that wrapper: nothing of his hand is on it — what
% shows is page 10's ink read through the sheet — and it is skipped, the gap in
% the page numbers being the record. Pages 2 to 8 are the first run, pages 10 to
% 15 the second; neither carries a pagination of his own, so both are followed
% by their headings. The inventory's title, « Cartes standards », is not the
% wrapper's: it comes from page 10, where he writes « cartes standard » in his
% own quotation marks, and it is therefore his too, developed by the
% archivists from the second run and not from the first.
%
% The first run is one proof. Page 2 sets up a functor Phi from the category of
% root systems to the category of systems (W, R, I, phi) — W a group, I a finite
% set, R a W-torsor, phi : R x I --> W — subject to four conditions: the
% equivariance phi(wr,i) = int(w) phi(r,i); that i |--> phi(r,i) be a system of
% generators making W a « groupe de Weyl épinglé »; the Coxeter presentation
% phi(r,i)^2 = 1, (phi(r,i)phi(r,j))^{n_ij} = 1; and that I carrying the n_ij be
% a Dynkin diagram. Its lower half is a heavily cancelled block on whether R may
% be read as the set of admissible generating systems of W parametrised by I —
% it may, exactly when the centre of W is trivial — and it closes by asking
% whether Phi is an equivalence: evidently faithful and essentially surjective,
% is it full? Pages 3 and 4 prove that it is, by showing Aut(Sigma) -->
% Aut(Phi(Sigma)) an isomorphism: alpha(w) = (int(w), w_R, id_I) is injective
% with image the automorphisms inducing the identity on I; a base point r_0
% turns such an automorphism into u_W = int(a), u_R = a_R, so the image is all
% of it; and a diagram of two short exact sequences over Aut(I) — which he
% accolades as « gr. des automorphismes ext. de Sigma » — closes the argument by
% the five lemma. The end of page 4 and page 5 change subject: the parabolic
% parts of Sigma correspond to the facettes of the hyperplane system in the dual
% V, the inclusion reversing, so that the strict flags form a simplicial object
% which is manifestly a triangulation of the sphere S^{ell-1}, and the set of
% repères is R x Rep'(I) with Rep'(I) = Bij([1,ell], I). Pages 6 to 8 are the
% computation that the run exists for: the operations sigma_0, ..., sigma_{ell-1}
% on R x Rep'(I) are the transpositions tau_{i+1} in the second factor except
% sigma_{ell-1}, which is (r,u) |--> (phi(r,u(ell)).r, u); they are involutions,
% they commute when j >= i+2, they satisfy the braid relation, and the order of
% pi = sigma_{ell-2}sigma_{ell-1} is computed by induction — pi^{2n}(r,u) =
% ((ab)^n r, u) with a = phi(r,u(ell-1)), b = phi(r,u(ell)) — to be 2nu, nu the
% lcm of the n_ij.
%
% The second run is a theorem in seven numbered items, stated and not proved:
% the numerals are circled on the leaves and are set here as (1) to (7), (3)
% being struck out on page 10 where a second circled numeral is blotted. (1)
% every non-empty connected and simply connected quasi-regular cellular map is
% regular and complete, with one exception, the plane containing a line, whose
% completion is the sphere with a punctured great circle. (2) such a map — the
% « cartes standard » — is determined up to isomorphism by the pair (p,q) of the
% order of its combinatorial vertices and of its faces, p and q running over
% N* u {infinity} with (1,p) and (p,1) excluded for p =/= 2; C_{p,q} with its
% repère r_{p,q} is the pinned standard map. (3) Aut(C_{p,q}) is the cartographic
% group modulo rho_s^p and rho_f^q, written Gamma_{p,q}. (4) the fundamental
% group of a quasi-regular map of type (p,q) embeds in Gamma_{p,q}, with
% equality in the regular pinned case. (5) brings in the Riemannian side: a
% « surface standard » is a non-empty connected simply connected real analytic
% Riemannian surface whose isometry group is transitive, and there are three of
% them up to a constant factor — the sphere under SO(3,R), the plane under the
% semi-direct product U.C of the Euclidean displacements, and the Poincaré
% half-plane under SL(2,R)/{±1} — named in the margins spherical, Euclidean and
% hyperbolic. A map is géodésique when its edges are arcs of geodesics,
% métriquement régulière when its isometric automorphisms act transitively on
% the repères, and géodésique standard when it is that and complete; a), b) and
% c) on pages 12 and 13 state existence, uniqueness of the isomorphism, and
% proper action, and the Corollaire states that the functor from standard
% geodesic normalised maps to standard cellular maps is an equivalence of
% categories. (6) sorts the C_{p,q}: spherical exactly for p = 2 or q = 2 and
% for (3,3), (3,4), (4,3), (3,5), (5,3) — the five Platonic solids, named on the
% leaf — equivalently when C_{p,q} is finite, equivalently when Gamma_{p,q} is
% finite; Euclidean exactly for (4,4), (3,6), (6,3), the three plane tilings,
% where Gamma_{p,q} is infinite and soluble; hyperbolic otherwise, where
% Gamma_{p,q} contains a free group on n generators for every n. (7) recovers
% the regular maps of type (p,q) from the subgroups pi of Gamma_{p,q} acting
% freely on X_{p,q}.
%
% The hand is drafting throughout and that is the folder's governing fact. The
% statements, the displayed formulas and the underlined words carry; a large
% part of the connective prose does not, and pages 12, 13 and 15 are the worst
% of it — page 15 in particular is a palimpsest whose left third is written in
% biais at some 45 degrees. Following folder 29 batches 8 and 9, folder 52 and
% folder 41, what the leaves will not support is left \ill{} rather than
% invented, and the apparatus is dense in consequence.
%
% Notation fixed once here rather than page by page. The reading hazard of the
% folder is that an opening parenthesis followed by his « r » runs together into
% a shape that reads as « w »: « phi(r, u(ell)) » looks like « phi(w, u(ell)) »
% everywhere on pages 6 and 7, and the two lines of page 7 where he sets
% phi(r, u(ell-1)) = a and phi(r, u(ell)) = b settle it at high resolution —
% every first argument of phi in this folder is an element of R. « ½ » is his
% abbreviation for « semi- »: the « ½ direct » of page 11 is the semi-direct
% product U.C. He tops the Sigma of a root system with a bar, inconstantly and
% within the same sentence, and the bar is not transcribed. « gpe », « s-gpe »,
% « comb. », « ppcm », « pt » are groupe, sous-groupe, combinatoire, plus petit
% commun multiple, point. N* with a bar over it is N* u {infinity}, as the
% « 1 <= p, q <= +infinity » of page 10 says. The circled numerals of the second
% run are set (1) to (7). His own underlining is set in bold.
%
% Readings flagged and not settled: the word after « le diagramme des 5 » on
% page 4; the noun after « correspondant à une un » on page 5, which the sense
% wants to be a simplicial complex and which the strokes will not give; the
% three words inserted above « provient d'un » on page 5; the sign written
% between Rep'(I) and Bij([1,ell], I) on that same page, drawn as a doubled
% arrow; the letter written twice on page 5 where V and the dual V had just been
% distinguished, and which is a V both times; the two words closing page 12,
% read « Donc pris »; and the proper name closing page 15, read « Bourguignon »
% and underlined by him.
%
% Pass: Opus 5 (claude-opus-5), 2026-09-13 — first pass, unchecked against the pages by a human.
\documentclass[11pt,a4paper]{article}
% Le préambule commun à toutes les transcriptions du fonds Grothendieck.
%
% Il tient en une idée : **l'incertitude se compose, elle ne se résout pas.**
% Une transcription qui lisse un mot douteux a détruit la seule chose pour
% laquelle on la faisait — un lecteur doit pouvoir savoir, à chaque endroit,
% ce qui était sur le feuillet et ce qui a été ajouté. Les six macros qui
% suivent sont donc l'appareil critique, et non de la décoration.
%
% Les mêmes commandes sont reconnues par `scripts/render.mjs`, qui produit la
% vue de lecture du site. Une macro ajoutée ici doit l'être aussi là-bas,
% faute de quoi le rendu échouera bruyamment — ce qui est voulu : un rendu
% silencieusement faux est pire qu'un rendu absent.

\ProvidesPackage{grothendieck}[2026/08/08 Transcriptions du fonds Grothendieck]

\usepackage{amsmath,amssymb,amsthm}
\usepackage{mathrsfs}
\usepackage{stmaryrd}
% Les diagrammes commutatifs sont partout dans ce fonds : c'est la langue
% même de la géométrie algébrique de Grothendieck, et une transcription qui
% les rendrait en prose ne transcrirait rien.
\usepackage{tikz-cd}
% Français par défaut — c'est la langue des feuillets — mais l'anglais est
% chargé aussi : la traduction et le résumé sont des documents anglais, et la
% typographie française (espaces avant les deux-points) y serait une faute.
% Ces documents basculent par \selectlanguage{english} en tête de corps.
\usepackage[english,french]{babel}
\usepackage{fontspec}
\usepackage{geometry}
\geometry{a4paper,margin=25mm,marginparwidth=28mm}
\usepackage{marginnote}
\usepackage{xcolor}
% ulem, not soul. `soul`'s \st and \ul are built on a character-by-character
% scan that XeLaTeX's font machinery breaks: the document fails with an
% undefined control sequence rather than degrading. ulem does the same two jobs
% and is Unicode-engine safe. `normalem` stops it hijacking \emph, which the
% transcriptions use for Grothendieck's own emphasis.
\usepackage[normalem]{ulem}
\usepackage{hyperref}
\hypersetup{colorlinks=true,linkcolor=black,urlcolor=black,pdfborder={0 0 0}}

% --- Métadonnées du lot -----------------------------------------------------
% Reprises telles quelles par le script de rendu, qui les affiche en tête de
% la vue de lecture. Elles fixent ce que le fichier prétend couvrir : sans
% elles, une transcription flotte sans point d'attache dans un fonds de seize
% mille feuillets.

\newcommand{\folder}[1]{\gdef\grothFolder{#1}}
\newcommand{\batch}[1]{\gdef\grothBatch{#1}}
\newcommand{\leaves}[2]{\gdef\grothFirst{#1}\gdef\grothLast{#2}}
\newcommand{\foldertitle}[1]{\gdef\grothFoldertitle{#1}}
\gdef\grothFolder{?}\gdef\grothBatch{?}\gdef\grothFirst{?}\gdef\grothLast{?}\gdef\grothFoldertitle{}

% --- Le résumé --------------------------------------------------------------
%
% Il ouvre la lecture modernisée, sous le titre « Résumé », et il est composé
% en petit. Ce n'est pas un ornement : c'est le seul passage du document qui
% ne soit pas des mathématiques, et un lecteur doit pouvoir le reconnaître
% avant de l'avoir lu — puis le sauter s'il connaît déjà le sujet, ou s'y
% arrêter s'il ne le connaît pas. Le filet qui le suit marque la reprise du
% corps.

\newenvironment{resume}
  {\small\setlength{\parskip}{0.45em}}
  {\par\medskip\hrule\medskip}

% --- Mots-clés ---------------------------------------------------------------
%
% En anglais, en clôture du résumé de la lecture modernisée : le vocabulaire
% moderne sous lequel chercher ce que les feuillets construisent. Le manifeste
% du site (`npm run manifest`) les extrait pour étiqueter la cote entière —
% cette ligne est la source unique des tags, il n'y a pas de fichier à côté.
\newcommand{\keywords}[1]{\par\smallskip\noindent
  {\footnotesize\textsc{Keywords} --- \textit{#1}}\par}

% --- L'appareil critique ----------------------------------------------------

% Le numéro de feuillet, en marge. C'est l'ancre qui permet de régler un
% désaccord sur une formule en regardant *un* feuillet plutôt qu'une chemise
% de six cents. Le site s'en sert aussi pour tourner les pages du fac-similé
% quand on fait défiler la transcription.
\newcommand{\leaf}[1]{%
  \marginnote{\footnotesize\color{gray!70}\textsf{#1}}%
  \label{leaf:#1}\ignorespaces}

% Illisible : on ne devine pas. Un mot inventé qui se lit comme les autres est
% le pire résultat possible.
\newcommand{\ill}{\textcolor{red!60!black}{[\,\ldots\,]}}

% Lecture douteuse : le mot est donné, et signalé comme incertain.
\newcommand{\uncertain}[1]{\uline{#1}}

% Ajout de l'éditeur — une lettre manquante, un mot sous-entendu.
\newcommand{\add}[1]{\textcolor{blue!50!black}{[#1]}}

% Biffé par Grothendieck lui-même. Ce qu'il a rayé fait partie du document :
% une rature dit souvent où la pensée a changé de direction.
\newcommand{\struck}[1]{\sout{#1}}

% Note du transcripteur, sur l'état matériel du feuillet.
\newcommand{\note}[1]{\par\smallskip{\small\color{gray!60!black}\textit{[#1]}}\par\smallskip}

% Note marginale de Grothendieck, distincte du corps du texte.
\newcommand{\marginal}[1]{\par\smallskip\noindent\hspace{1em}%
  {\small\color{gray!60!black}#1}\par\smallskip}

% --- Titre ------------------------------------------------------------------

\AtBeginDocument{%
  \begin{center}
    {\large\bfseries Fonds Alexandre Grothendieck --- cote n\textsuperscript{o}~\grothFolder}\\[2pt]
    {\itshape\grothFoldertitle}\\[4pt]
    {\small feuillets \grothFirst--\grothLast\ (lot \grothBatch)}\\[2pt]
    {\footnotesize\color{gray!60!black}Transcription établie d'après le fac-similé numérisé
     par l'Université de Montpellier.\\
     Les lectures incertaines sont soulignées, les illisibles notées [\,\ldots\,],
     les ajouts entre crochets.}
  \end{center}
  \vspace{1em}}

\endinput


\folder{88}
\batch{1}
\pages{1}{15}
\dating{s.d. — le groupe « Géométrie et topologie combinatoire » (69 à 102) est daté [1976-vers 1986]}
\watermark{Édition de démonstration}
\foldertitle{Cartes standards : notes manuscrites (s.d.)}

\begin{document}

\section{Cartes associées aux syst. de Racines}

\note{ce titre est de sa main : il occupe seul, au crayon, le coin supérieur
droit du feuillet de couverture, qui ne porte rien d'autre et ne reçoit donc
pas de marqueur de page. Il est suivi d'un « (7) » dont le sens n'est pas
établi. Le titre de l'inventaire, « Cartes standards », ne vient pas de là mais
de la page 10}

\page{2}

Cat. des systèmes de racines $\overset{\Phi}{\longrightarrow}$ cat. des
systèmes
\[
  \bigl(W,\ R,\ I,\ \varphi : R \times I \longrightarrow W\bigr)
\]
où $W$ est un groupe ($I$ ens. fini), $R$ un $W$-torseur, $\varphi$ une
\uncertain{application}\note{le mot est en partie recouvert ; un $I$ et un $W$
biffés le précèdent sur la même ligne}, avec

\begin{itemize}
  \item[a)] $\varphi(wr, i) = \operatorname{int}(w)\,\varphi(r,i)$
  \item[b)] $\forall\, r \in R$, $i \mapsto \varphi(r,i)$ définit un système de
    générateurs de $W$ qui en fait un « groupe de Weyl épinglé », i.e.
  \item[c)] $\varphi(r,i)^2 = 1$, \quad
    $\bigl(\varphi(r,i)\,\varphi(r,j)\bigr)^{n_{ij}} = 1$ \ \ $(i \neq j)$

    $\longrightarrow$ une présentation de $W$, où $n_{ij}$
    \uncertain{est l'ordre} de $\varphi(r,i)\,\varphi(r,j)$, et
  \item[d)] $I$ muni du syst. des $n_{ij}$ ($n_{ij} = n_{ji}$) est un
    diagramme de Dynkin ($i$ lié à $j$ ssi $n_{ij} \geqslant 3$, i.e.
    $n_{ij} \neq 2$, i.e. $\varphi(r,i)$ ne commutant pas à $\varphi(r,j)$).
\end{itemize}

\textbf{NB} On a
\[
  R \overset{\Psi}{\longrightarrow} \text{Syst. gén. admissibles de } W
  \text{ paramétrés par le diagramme de Dynkin } I
\]

Mais cette identification, \ill{} compatible avec l'action de $W$, n'est pas
\add{en gén.} injective \ill{}

\note{suit un bloc d'une dizaine de lignes annulé par de longs traits
horizontaux et deux diagonales, écrit par-dessus une première version
elle-même biffée. Ce qui s'en lit n'excède pas : « (si $I \neq \emptyset$),
$r \in R$ », « $r'$ sont deux repères », « chose équivalent : deux syst. de
repères », « l'élément (dans groupe de $W$) qui envoie \ill{} le stabilisateur
de repères, dans $W$ ». Le bloc n'est pas donné}

\struck{\ill{}} l'élément neutre. Or ce stabilisateur est évidemment le
\textbf{centre} de $W$ — donc $\Psi$ injection (permet d'interpréter $R$ comme
un ens. des syst. de gén. admissibles de $W$ paramétrés par le diagramme de
Dynkin $I$) ssi $Z(W) = \{1\}$.

Le foncteur précédent est-il une équivalence de catégories ? Il est évid.
fid., \uncertain{ess.} surjectif, est-il pl. fidèle ? Deux syst. de racines
$\Sigma$, $\Sigma'$ tels que $\Phi(\Sigma) \simeq \Phi(\Sigma')$ ont-ils des
diagrammes de Dynkin isomorphes,

\page{3}

donc sont isomorphes. Il reste à voir que pour tt $\Sigma$,
\[
  \operatorname{Aut}(\Sigma) \overset{\sim}{\longrightarrow}
  \operatorname{Aut}(\Phi(\Sigma))
\]
est un iso de gpes. Or on a un hom. Considérons
\[
  W \overset{\alpha}{\longrightarrow} \operatorname{Aut}(\Phi(\Sigma))
\]
par
\[
  \alpha(w) = \bigl(\operatorname{int}(w),\ w_R,\ \mathrm{id}_I\bigr)
\]

en effet,
\[
  w_R(w'r) = \operatorname{int}(w)(w')\ w_R\,r
\]
\note{les deux membres sont accolés au-dessous, à gauche $w(w'r)$, à droite
$(w w' w^{-1})(wr)$ ; les lettres $w' \in W$ et $r \in R$ sont rappelées
au-dessus de la ligne}
et
\[
  \varphi\bigl(w_R(r),\ i\bigr) = \operatorname{int}(w)\,\varphi(r,i)
\]
\note{le membre de gauche est repris au-dessous sous la forme
$\varphi(wr,i)$ ; les deux lignes sont accolées ensemble « pour a) »}

Ceci dit, $\alpha$ est \add{$(u_W, u_R)$} évidemment injectif, son image est
formée des automorphismes de $\Phi(\Sigma) = (W,R,I,\varphi)$ induisant
l'identité sur $I$, i.e. de la forme $(u_W, u_R, \mathrm{id}_I)$, où $u_R$ est
un automorphisme affine de $R$, associé à l'automorphisme $u_W$ du gpe $W$
(\uncertain{$R$} « partie homogène » de $u_R$)\note{le sens demande ici $u_W$ ;
le feuillet porte un $R$} et satisfaisant à
\[
  (\ast\ast) \qquad \varphi(u_R\,r,\ i) = u_W\,\varphi(r,i)
  \qquad \text{i.e.} \qquad
  \varphi(r,i) = u_W^{-1}\,\varphi(u_R\,r,\ i)
\]

Choisissons un point $r_0 \in R$ [\struck{\ill{}} d'identifier $R : W$] donc
\ill{}
\[
  u_R(w\,r_0) = u_W(w)\,a\,r_0
  \qquad \text{où } \struck{\ill{}}\ u_R(r_0) = a\,r_0
\]
et la condition $(\ast\ast)$ s'écrit (pour $r = w\,r_0$ \ill{})

\page{4}

\[
  \varphi\bigl(u_W(w)\,a\,r_0,\ i\bigr) = u_W\bigl(\varphi(w\,r_0,\ i)\bigr)
\]

\note{le développement est disposé en deux colonnes reliées par des
« $\parallel$ ». À gauche $\operatorname{int}(u(w))\,\varphi(a\,r_0, i)$, puis
$u(w)\,a\,\varphi(r_0,i)\,a^{-1}\,u(w)^{-1}$ ; à droite
$u\bigl(w\,\varphi(r_0,i)\,w^{-1}\bigr)$, puis
$u(w)\,u(\varphi(r_0,i))\,u(w)^{-1}$. Une première version du membre de gauche
est biffée}

on a donc
\[
  a\,\varphi(r_0,i)\,a^{-1} = u\,\varphi(r_0,i)
  \qquad \forall\, i \in I
\]
\note{« $i \in S$ » est inséré au-dessus de « $\forall\, i \in I$ »}
i.e. $u_W = \operatorname{int}(a)$ (car les $\varphi(r_0,i)$
\uncertain{engendrent} $W$)\note{le mot est écrit par-dessus un autre, biffé}
donc
\[
  u_R(w\,r_0) = (a\,w\,a^{-1})(a\,r_0) = a(w\,r_0) = a_R(w\,r_0)
\]
i.e. $u_R = a_R$, i.e. $(u_W,\ \struck{\ill{}}\ u_R,\ \mathrm{id}_I)
= \alpha(a)$.

On a \struck{\ill{}} un diagramme de suites exactes

\begin{tikzcd}[column sep=small, row sep=small, nodes={font=\scriptsize}]
  1 \arrow[r] & W \arrow[r, "\alpha_0"] \arrow[d, "\mathrm{id}"] &
    \operatorname{Aut}(\Sigma) \arrow[r, "p_0"] \arrow[d] &
    \operatorname{Aut}(I) \arrow[r] \arrow[d, "\mathrm{id}"] & 1 \\
  1 \arrow[r] & W \arrow[r, "\alpha"'] &
    \operatorname{Aut}(\Phi(\Sigma)) \arrow[r, "p"'] &
    \operatorname{Aut}(I) &
\end{tikzcd}

\marginal{gr. des automorphismes ext. de $\Sigma$}

\note{cette dernière ligne est de lui, écrite au-dessus du $\operatorname{Aut}(I)$
de la ligne supérieure et accolée à lui par un « $\parallel$ »}

et le diagramme des 5 \ill{}, montre alors que la flèche verticale médiane est
un iso, cqfd.

Les \struck{\ill{}} parties paraboliques \uncertain{des racines} de
$\Sigma \subset V$ \ill{} \struck{correspondent aux} \ill{} aux facettes
\add{fermées} du syst. des hyperplans dans $\check{V}$ associé au syst. des
racines dans $V$.

\note{le passage est repris trois fois : deux débuts biffés, une insertion
interlinéaire entre parenthèses que le feuillet ne soutient pas, et une
seconde insertion au-dessus de « facettes »}

\page{5}

l'inclusion des parties paraboliques (ou encore, des gpes paraboliques
correspondants, quand le syst. de racines $\Sigma$ provient d'un \ill{} groupe
alg.)\note{trois mots sont insérés au-dessus de « provient d'un » et ne se
lisent pas} \struck{est} se transforme par \ill{} l'inclusion des facettes dans
$\check{V}$. C'est cette relation (opposée à l'inclusion des paraboliques)
qu'on prend sur cet \struck{\ill{}} ens. de parties, qui devient dès lors une
géométrie des \textbf{drapeaux stricts}, correspondant à une un \ill{}
simplicial\note{« une » reste en fin de ligne, non biffé, et le nom qui suit ne
se lit pas}, dont $V$ sont des sommets (correspond aux « paraboliques
minimaux » \add{ou drapeaux minimaux}), et $V$ des simplexes maximaux aux
« Borels » ou « syst. simples de racines » ou repères $r \in R$.

\note{la même lettre est écrite aux deux endroits, alors que la ligne
précédente distingue $V$ de $\check{V}$}

En fait, il s'agit manifestement d'une triangulation de la sphère
$S^{\ell-1}$. L'ens. des repères s'identifie à
$R \times \struck{\ill{}}\ \mathrm{Rep}'(I)$, où
\[
  \mathrm{Rep}'(I) \Longrightarrow \mathrm{Bij}\bigl([1,\ell],\ I\bigr)
  \longrightarrow \text{l'ens. des ordres totaux sur l'ens. } I
  \text{ (de card. } \ell) .
\]
\note{le premier signe est tracé comme une double flèche et vaut ici une
identification}

\page{6}

\ill{} est un torseur (à droite) sous $\mathfrak{S}_\ell$. Les op.
$\sigma_0, \ldots, \sigma_\ell$ sur $R \times \mathrm{Rep}'(I)$ s'identifient
\ill{} ainsi : $\sigma_0, \sigma_1, \ldots, \sigma_{\ell-1}$ sont les
opérations $\mathrm{id}_R \times \ill{}\ \tau_{i+1}$
($0 \leqslant i \leqslant \ell-2$), où $\tau_{i+1} \in \mathfrak{S}_\ell$
\ill{} la transposition de $i+1$ et $i+2$ ;
\[
  \sigma_{i'}\bigl(r,\ u : [1,\ell] \overset{\sim}{\to} I\bigr)
  = (r,\ u \circ \tau_{i'}),
  \qquad 0 \leqslant i' \leqslant \ell-1
\]
\[
  \sigma_{\ell-1}\bigl(r,\ u : [1,\ell] \longrightarrow I\bigr)
  = \bigl(\varphi(r,\ u(\ell))\cdot r,\ u\bigr)
\]
\note{une insertion interlinéaire court au-dessus de la première de ces deux
lignes ; seul « à $\sigma_{\ell-1}$ » s'y lit. Les deux bornes se recoupent :
la seconde ligne redéfinit le $\sigma_{\ell-1}$ que la première venait
d'inclure}

On trouve bien que $\sigma_i^2 = 1$ ($0 \leqslant i \leqslant \ell-1$) et que
$\sigma_i, \sigma_j$ commutent si $j \geqslant i+2$
($0 \leqslant i, j \leqslant \ell-1$) — le seul cas non trivial étant
$j = \ell-1$, $i \leqslant \ell-3$,
\[
  \sigma_i\,\sigma_{\ell-1}(r,u)
  = \sigma_i\bigl(\varphi(r,u(\ell))\cdot r,\ u\bigr)
  = \bigl(\varphi(r,u(\ell))\cdot r,\ u \circ \tau_{i+1}\bigr)
\]
\[
  \sigma_{\ell-1}\,\sigma_i(r,u)
  = \sigma_{\ell-1}\bigl(r,\ u \circ \tau_{i+1}\bigr)
  = \bigl(\varphi\bigl(r,\ (u \circ \tau_{i+1})(\ell)\bigr)\cdot r,\
    u \circ \tau_{i+1}\bigr)
\]
\note{$(u \circ \tau_{i+1})(\ell)$ est accolé à $u(\ell)$, avec
« car $i \leqslant \ell-3$ i.e. $i+1 \leqslant \ell-2$ »}

Ici $(\sigma_i\,\sigma_{i+1})^3 = \mathrm{id}$ — distinguer deux cas :
$0 \leqslant i \leqslant \ell-2$, \ill{} cela résulte de
$(\tau_{i+1}\,\tau_{i+2})^3 = 1$ ; puis $i = \ell-1$, i.e.
$\sigma_{\ell-1}\,\sigma_{\ell-2}$ — \ill{}
\[
  \sigma_{\ell-1}\,\sigma_{\ell-2}(r,u)
  = \sigma_{\ell-1}\bigl(r,\ u \circ \tau_{\ell-1}\bigr)
  = \bigl(\varphi\bigl(r,\ (u \circ \tau_{\ell-1})(\ell)\bigr)\cdot r,\
    u \circ \tau_{\ell-1}\bigr)
\]
\note{le second argument de $\varphi$ est accolé à $u(\ell-1)$}
\[
  \pi(r,u) = \bigl(\varphi(r,\ u(\ell-1))\cdot r,\ u \circ \tau_{\ell-1}\bigr)
\]
\note{$\pi$ n'est introduit par aucune définition ; c'est
$\sigma_{\ell-1}\sigma_{\ell-2}$, ce que la page 8 confirme. Il note $w$
au-dessus de $\varphi(r,u(\ell-1))$, et au-dessous $r'$ sous
$\varphi(r,u(\ell-1))\cdot r$ et $u'$ sous $u \circ \tau_{\ell-1}$}
\[
  \pi^2(r,u) = \struck{\ill{}}\
  \bigl(\varphi(r',\ u'(\ell-1))\cdot r',\ u' \circ \tau_{\ell-1}\bigr)
\]
\note{$u'(\ell-1)$ est accolé à $u(\ell)$ et $u' \circ \tau_{\ell-1}$ à $u$ ;
une première version de la ligne, $\varphi(\varphi(r,u(\ell-1))r,\ u)$, est
biffée}
\[
  = \bigl(\varphi(w\cdot r,\ u(\ell))\cdot w\,r,\ u\bigr)
\]

\page{7}

\[
  \pi^2(r,u)
  = \bigl(w\,\varphi(r,u(\ell))\,w^{-1}\,w\,r,\ u\bigr)
  = \bigl(w\,\varphi(r,u(\ell))\,r,\ u\bigr)
\]
\[
  = \bigl(\varphi(r,u(\ell-1))\,\varphi(r,u(\ell))\,r,\ u\bigr)
\]
\[
  \pi^3(r,u)
  = \bigl(\varphi(r'',\ u(\ell-1))\,r'',\ u \circ \tau_{\ell-1}\bigr)
  \quad \struck{cqf}
\]
\note{$r'' = w'\,r$ ; l'accolade développe
$w'\,\varphi(r,u(\ell-1))\,w'^{-1}\,r'' = w'\,\varphi(r,u(\ell-1))\cdot r$}

\struck{\ill{}}
\[
  \pi^3(r,u) = \bigl(\varphi(r,u(\ell-1))\,\varphi(r,u(\ell))\,
    \varphi(r,u(\ell-1))\cdot r,\ u \circ \tau_{\ell-1}\bigr)
\]
\note{deux amorces sont biffées au-dessus de cette ligne, dont une
$\pi^{2n}(r,u) =$}
\[
  \pi^4(r,u)
  = \bigl(\varphi(r''',\ u'(\ell-1))\,r''',\ u' \circ \tau_{\ell-1}\bigr)
  = \bigl(w'''\,\varphi(r,u(\ell))\,r,\ u\bigr)
  = \Bigl[\bigl(\varphi(r,u(\ell-1))\,\varphi(r,u(\ell))\bigr)^2 r,\ u\Bigr]
\]

Soit
\[
  \varphi(r,\ u(\ell-1)) = a, \qquad \varphi(r,\ u(\ell)) = b .
\]

On trouve par récurrence
\[
  \begin{cases}
    \pi^{2n}(r,u) = \bigl((ab)^n\,r,\ u\bigr) \\[2pt]
    \pi^{2n+1}(r,u) = \bigl(((ab)^n a)\,r,\ u \circ \tau_{\ell-1}\bigr)
  \end{cases}
\]
\note{$r_n = w_n\,r$ est accolé sous la première ligne, $r'_n = w'_n\,r$ et
$u'$ sous la seconde}

Donc
\[
  \pi^i(r,u) = (r,u)
  \iff
  \begin{cases}
    i \equiv 0 \ (2), \text{ i.e. } i = 2n \\[2pt]
    (ab)^n = 1
  \end{cases}
\]
i.e. $n_{u(\ell-1),\,u(\ell)} \mid n$, i.e. \struck{$i \equiv 0\ (2n$}
\uncertain{$n$} $\equiv 0 \ (2\,n_{u(\ell-1),\,u(\ell)})$

\note{le « $=\ \mathrm{id}$ » qui suivait $\pi^i$ est biffé}

Suit, entre crochets, le pas de la récurrence :
\[
  \pi^{2n+1}(r,u)
  = \bigl(\varphi(r'_n,\ u'(\ell-1))\cdot r'_n,\ u' \circ \tau_{\ell-1}\bigr)
  = \struck{\ill{}}\ \bigl((ab)^{n+1}\,r,\ u\bigr)
\]
\[
  \pi^{2n+2}(r,u)
  = \bigl(\varphi(r_{n+1},\ u(\ell-1))\cdot r_{n+1},\
    u \circ \tau_{\ell-1}\bigr)
\]
\note{accolades : $w'_n\,\varphi(r,u(\ell))\,w'^{-1}_n\,w'_n\,r$, avec
$(ab)^n a$ sous $w'_n$ et $b$ sous $\varphi(r,u(\ell))$ ; puis
$w_{n+1}\,\varphi(r,u(\ell-1))\,w_{n+1}^{-1}\,w_{n+1}\,r = (ab)^{n+1} a\,r$}

\page{8}

\[
  \pi^i(r,u) = (r,u) \quad \forall\, (r,u) \in R \times \mathrm{Rep}(I)
  \iff
  \begin{cases}
    i \equiv 0\ (2), \text{ i.e. } i = 2n \\[2pt]
    \bigl(\varphi(r,u(\ell-1))\,\varphi(r,u(\ell))\bigr)^n = 1
    \quad \forall\, r, u
  \end{cases}
\]
\note{un « i.e. » suivi de trois ou quatre mots est biffé au bout de la
seconde accolade}

i.e. $n$ multiple \uncertain{des} ppcm des $n_{ij}$
($i, j \in I$, $i \neq j$)

Donc on trouve
\[
  \begin{cases}
    \text{ordre de } \pi = \sigma_{\ell-2}\,\sigma_{\ell-1} = 2\nu \\[2pt]
    \text{où } \nu = \text{le ppcm des } n_{ij}
      \ (i, j \in I,\ i \neq j)
  \end{cases}
\]
\note{sic : l'ordre de $\pi$ est égalé à $\pi$ lui-même autant qu'à $2\nu$}

\section{Théorème (à prouver)}

\note{ce titre est de sa main, en tête de la page 10, et ouvre le second
feuillet du dossier. Les numéros d'articles sont cerclés sur le feuillet et
sont donnés ici sous la forme (1), (2), \ldots}

\page{10}

\begin{itemize}
  \item[(1)] Toute carte cellulaire \add{$\neq \emptyset$} connexe et
    simplement connexe quasi régulière (i.e. dont tous les sommets comb. et
    toutes les faces ont même ordre) est régulière, et elle est complète (i.e.
    les sommets comb. d'ordre fini sont effectifs) sauf dans le cas d'une carte
    isomorphe : $(\mathbb{R}^2 \supset \mathbb{R})$ (dont la complétion \ill{}
    la sphère munie d'un grand-cercle ponctué\ldots).

  \item[(2)] Une carte cellulaire \add{$\neq \emptyset$} conn. simpl. connexe
    régulière \add{une telle carte} et complète (\struck{\ill{}}
    « cartes standard ») est déterminée, à isomorphisme près, par
    \add{le couple $(p,q)$} l'ordre $p$ des sommets comb., $q$ l'ordre des
    faces ($p, q \in \overline{\mathbb{N}}^*$, entiers,
    $1 \leqslant p, q \leqslant +\infty$). Le couple $(p,q)$ peut prendre
    toutes les valeurs \struck{\ill{}} $\in \overline{\mathbb{N}}^* \times
    \overline{\mathbb{N}}^*$ sauf $(1,p)$ \add{et $(p,1)$ si $p \neq 2$}.
    [Appelons $(C_{p,q},\ r_{p,q})$ la carte standard épinglée (i.e. munie d'un
    bi-drapeau repère $r_{p,q}$) correspondant à $(q,p)$, déterminée à
    \struck{homéomorphisme près, unique} isomorphisme près, pris dans la
    catégorie \uncertain{isotopique} des cartes.]

  \item[(3)] \struck{\ill{}}
    \[
      \operatorname{Aut}(C_{p,q})
      \overset{\sim}{\longleftarrow}
      \widehat{\Gamma}\big/\bigl(\rho_s^{\,p},\ \rho_f^{\,q}\bigr)
      \overset{\text{déf}}{=} \Gamma_{p,q}
    \]
    \note{le $\widehat{\Gamma}$ est fortement encré et porte au-dessous
    « \uncertain{gpe} cartographique » ; le $\Gamma$ de $\Gamma_{p,q}$ est
    encré de même. Un second numéral cerclé, entre (3) et (4), est entièrement
    noirci}

  \item[(4)] \textbf{Corollaire} Soit $\mathcal{C} = (X \supset K \supset S)$
    une carte cellulaire complète \ill{} connexe quasi-régulière de type
    $(p,q)$. Alors le revêtement universel $\widetilde{\mathcal{C}}$
    correspondant, le groupe fondamental
\end{itemize}

\page{11}

\ldots\ tal de $X$ s'identifie : \ill{} un sous-groupe de \struck{$\Gamma$}
$\Gamma_{p,q}$. Si $\mathcal{C}$ est régulière \add{épinglée}, \ill{} le gpe
des \struck{\ill{}} automorphismes, dans le gpe fondamental — \ill{} un
\uncertain{modèle} mixte de $(X, G)$ s'identifie à $\Gamma_{p,q}$.

\begin{itemize}
  \item[(5)] [Appelons \struck{\ill{}} « surface (riemannienne) standard » une
    surface \struck{$E$} analytique réelle munie d'une métrique
    \add{riemannienne $\geqslant 0$} \struck{telle que} :
\end{itemize}

\begin{itemize}
  \item[a)] $E$ est $\neq \emptyset$, connexe, simplement connexe
  \item[b)] Le groupe des isométries de $E$ est transitif sur $E$.
\end{itemize}

On sait alors que, \ill{}, à un facteur constant près sur la métrique, $E$ est
isomorphe à l'un des trois espaces \ill{} :

\begin{itemize}
  \item[1\textsuperscript{o})] La sphère $S^2$, espace homogène sous
    $SO(3,\mathbb{R})$ \add{avec $\{z, \bar{z}\}$} \ill{}
  \item[2\textsuperscript{o})] Le plan \struck{euclidien} \add{complexe}, sous
    le groupe $\mathcal{U}\cdot\mathbb{C}$ (½ direct) le gpe des déplacements
    euclidiens
  \item[3\textsuperscript{o})] Le demi-plan de Poincaré,
    \struck{comme espace} espace homogène de $SL(2,\mathbb{R})/\{\pm 1\}$.
\end{itemize}

\marginal{— plan sphérique \ill{} déplacements sphériques}

\marginal{plan \ill{}, \ill{} euclidiens}

\marginal{plan hyperbolique, dépl. hyperboliques}

\note{ces trois notes sont écrites en biais, à quelque 45°, dans la marge de
gauche, en regard des trois cas}

Les gpes écrits sont les comp. \struck{connexes} \add{fermées} des groupes des
isométries — ou mieux, des isométries qui conservent l'orientation — ou encore,
les \ldots

\page{12}

\ldots\ structure complexe associée à l'orientation (dans le cas où une telle
orientation a été choisie). Notez que dans les cas 1\textsuperscript{o} et
3\textsuperscript{o}, \struck{\ill{}} les similitudes sont des isométries, on
peut normaliser le \struck{métrique} — en particulier dans la « surface
standard normalisée » — déterminée pour chacun des \struck{trois} trois types :
\uncertain{Donc pris}.

Une carte cellulaire sur une surface riemannienne $X$ est dite
\textbf{géodésique} si \struck{toutes} ses arêtes sont des \add{arcs de}
\textbf{géodésiques} ; elle est dite \textbf{métriquement régulière} si elle
est \textbf{géodésique} et si le gpe des automorphismes \textbf{isométriques}
de la carte est transitif sur l'un des repères (ou simplement transitif sur
\ill{} des repères).

On dit qu'elle est une \textbf{carte géodésique standard} si elle est
\textbf{métriquement} régulière \add{complète (cf. 1\textsuperscript{o}))},
\struck{et si} \ill{} \struck{\ill{} est une surface riemannienne
\ill{} standard \ill{}} \struck{Ceci} \add{Donc on doit avoir ceci}
\struck{On dit qu'elle est de type}

\begin{itemize}
  \item[a)] Pour toute carte standard, il existe \struck{une} une carte
    géodésique standard $(X_0, K_0, S_0)$ qui lui est
    \uncertain{topologiquement} isomorphe.
  \item[b)] \struck{\ill{}} Pour deux cartes géodésiques standard, dans la
    classe d'isotopie d'isomorphismes \ldots
\end{itemize}

\note{a) et b) sont enfermés dans un double filet vertical tracé dans la marge}

\page{13}

\ldots\ topologiques il y a un isomorphisme topologique unique qui soit une
\struck{\ill{}} \textbf{isométrie} (dans une isométrie dans le cas où les
surfaces ne sont pas \struck{hyperboliques} \textbf{sphériques}, \ill{}
hyperboliques \ill{}) unique

\begin{itemize}
  \item[c)] Le groupe des automorphismes \ill{} d'une carte géod. st.
    $(X, K, S)$ opère \textbf{proprement} sur $X$ \ill{}
\end{itemize}

\textbf{NB} On dit que la carte \ill{} est \textbf{normalisée} si \ill{}
\ill{}\ hyperbolique, $(X, g)$ \ill{} le cas euclidien \ill{} carte
\struck{\ill{}} \textbf{loc. f. finie} \ill{} des autres \ill{} \ill{} que les
longueurs \ill{} finies dans les \textbf{longueurs} \ill{} des arêtes soit
égale \uncertain{à} 1 (on \ill{} \struck{à la plus} exige que \ill{} égale :
1\ldots). D\ill{} \ldots

\note{tout ce passage est écrit vite, sur trois niveaux d'insertions, et la
prose de liaison n'y est pas recouvrable ; ce qui suit reprend au
\textbf{Corollaire}}

\textbf{Corollaire} \ldots\ \ill{} foncteur « \uncertain{métrique} » de la
catégorie des cartes \textbf{géodésiques standard normalisées} (avec comme
morphismes les isométries qui respectent la carte) dans la catégorie des cartes
\textbf{cellulaires standard}, est une \textbf{équivalence de catégories}.

\textbf{R. NB} Carte géodésique standard normalisée \textbf{épinglée} : les
surfaces ambiantes \uncertain{s'identifient}.

\page{14}

\ldots\ canoniquement aux surfaces riemanniennes standard \struck{type}
$S^2$, $\mathbb{R}^2 = \mathbb{C}$, $\mathbb{H}$ avec le « repère »
\ill{} habituel \ldots\ On \ill{} encore noter les $C_{p,q}$.

\struck{\ill{}} On \textbf{distingue} \struck{types} des cartes standard de
type sphérique, euclidien \ldots\ ou hyperbolique.

\begin{itemize}
  \item[(6)] \struck{\ill{}} $C_{p,q}$ est \uncertain{sphérique}
    \struck{\ill{}} exactement dans les cas suivants :
    \struck{$p = q$} $p = 2$ ($1 \leqslant q < +\infty$) ; $q = 2$
    ($1 \leqslant p < +\infty$), \struck{$p = q = 3$},
    \[
      (3,3) \quad (3,4) \quad (4,3) \quad (3,5) \quad (5,3)
    \]
    \[
      \text{tétraèdre} \quad \text{cube} \quad \text{octaèdre} \quad
      \text{dodécaèdre} \quad \text{icosaèdre}
    \]
    — ou encore ssi $C_{p,q}$ \textbf{finie} ou encore ssi $\Gamma_{p,q}$ fini.

  \item[(b)] $C_{p,q}$ est euclidien exactement dans les cas suivants :
    \struck{$(2,\infty)$} \struck{$(\infty,2)$}
    \[
      (4,4), \quad (3,6), \quad (6,3)
    \]
    \[
      \text{pavage carré} \quad \text{pavage \ill{}} \quad
      \text{pavage triangulaire}
    \]
    ses $\Gamma_{p,q}$ est infini \ill{} \struck{ni} et \textbf{résoluble}.

  \item[(c)] $C_{p,q}$ est hyperbolique dans les autres cas, i.e.
    \struck{\ill{} $(3,3)$} $(p \geqslant 3,\ \ill{}\ q \geqslant 3)$ ou
    $(p = 3,\ \ill{})$
\end{itemize}

\begin{itemize}
  \item[$(c_1)$] ou \struck{\ill{}} $(p,q) > (4,4)$ ; ou $p = 3$,
    $q \geqslant 7$ ; ou $q = 3$, $p \geqslant 7$ ; ou $(p,q) = (4,5)$ ou
    $(p,q) = (5,4)$
  \item[$(c_2)$] ou encore \struck{\ill{}} $p + q \geqslant 10$ ou
    $(p,q) \neq (3,6), (6,3)$, \ill{}
  \item[$(c'_2)$] ou encore $\Gamma_{p,q}$ contient un sous-groupe
    \textbf{libre} à $2$ générateurs \ill{} libre à $n$ générateurs
    ($\forall\, n \neq 0$)
\end{itemize}

\note{$(c_1)$ est accolé sur trois lignes du feuillet, ici mises bout à bout}

(dans ce cas \struck{a-t-on} \ill{} \uncertain{quotient} \ill{})

\begin{itemize}
  \item[(7)] Les cartes régulières $(X \supset K \supset S)$ de type $(p,q)$,
    \ill{} épinglées, s'obtiennent \ill{} canoniquement \ill{} pris dans la
    cat. isotopique des cartes) où à partir des \ill{} des sous-groupes
    $\pi \subset \Gamma_{p,q}$
\end{itemize}

\page{15}

\ldots\ opérant \struck{\ill{}} \add{librement} sur $X_{p,q}$, \ill{} (qui
s'identifie au quotient de $C_{p,q}$ par $\pi$, $X$\ldots) Telles que le
\textbf{groupe fondamental} de \ill{} \ill{} s'identifie d'une manière
canonique — \ill{} isotopie de la carte) d'une \ill{} \textbf{métrique
riemannienne} (la métrique \ill{} \textbf{géodésique}) \ill{} pour laquelle la
carte est \textbf{régulière} \ldots\ les $\pi$ de \ill{} $\Gamma_{p,q}$ \ill{}
$\Gamma_{p,q}$ est \textbf{invariant}.

\textbf{NB} $\pi$ \ill{} des liens, pour \ill{} $(p,q)$ donné ; pour
déterminer les $\pi$ possibles, d'étudier \ill{} $\Gamma_{p,q}$ dans
$X_{p,q}$ — et plus \ill{} les \struck{\ill{}} \ill{} stabilisateurs des
$\Gamma_{p,q}$ des pts de $X_{p,q}$.

\struck{\ill{}} \textbf{Re NB} Il y a la \ill{} \ill{} métrique \ill{} \ill{}
\uncertain{correspondant} \ill{} (i.e. \struck{\ill{}} fini et $\pi$ d'indice
fini) \ill{} \textbf{complexes} — qui, d\ill{} fait dans une \ill{} seule !

La carte (\struck{n. rig.}) est orientée (p. ex. orientable de \ill{}
\struck{\ill{}} épinglée) \ill{} (partie régulière de \ill{}) d'une structure
(i.e. carte \textbf{finie}) \ill{}, dans le cas compact \ill{} correspondre
\ill{} algébriques projectives \ill{} \struck{\ill{}}
\uncertain{Bourguignon}.

\note{la dernière page est un palimpseste : son tiers gauche est écrit en biais
à quelque 45° et n'est pas recouvrable, et la prose de liaison de la colonne de
droite ne l'est qu'en partie. Le nom qui la clôt est souligné par lui}

\end{document}
